\section{The G\"odel sentence}
\label{sec:first-incompleteness-03}

Let \(T\) be a computably enumerable theory extending \(\mathsf Q\).  Apply
\cref{thm:diagonal-lemma}, in the complexity-preserving form discussed
below display (G), to
\(\theta(x)=\neg\mathsf{Prov}_T(x)\).  Obtain a \(\Pi^0_1\) sentence \(G_T\)
satisfying
\[
  T\derives G_T\leftrightarrow
  \neg\mathsf{Prov}_T(\overline{\ulcorner G_T\urcorner}).
  \tag{G}
\]
The fixed-point biconditional is in fact derivable already in \(\mathsf Q\).
The complexity-preservation clause is a standard refinement of the coding
construction: choose a suitable representation of substitution on codes and
put the resulting fixed point back into prenex form.  The fixed point of a
\(\Pi^0_1\) formula may then be chosen \(\Pi^0_1\).  This clause is cited
rather than proved here; see
\cite{Shoenfield1967}.

\begin{theorem}[G\"odel's first incompleteness theorem, original form]
\label{thm:godel-first}
Let \(T\) be a computably enumerable theory extending \(\mathsf Q\).
\begin{enumerate}
  \item If \(T\) is consistent, then \(T\notderives G_T\).
  \item If \(T\) is \(\Sigma^0_1\)-sound, then
        \(T\notderives\neg G_T\).
  \item Omega-consistency also suffices for
        \(T\notderives\neg G_T\).
\end{enumerate}
Consequently, if \(T\) is omega-consistent or \(\Sigma^0_1\)-sound, it is not
deductively complete.
\textit{Status: proof sketch.}
\end{theorem}

\begin{proofplan}
Suppose \(T\derives G_T\).  Because a concrete \(T\)-proof exists and the
proof predicate represents concrete proofs,
\(T\derives\mathsf{Prov}_T(\overline{\ulcorner G_T\urcorner})\).  But (G)
also gives
\(T\derives\neg\mathsf{Prov}_T(\overline{\ulcorner G_T\urcorner})\).
Thus \(T\) is inconsistent.  This proves (1).

If \(T\) is consistent, (1) says that no standard number codes a
\(T\)-proof of \(G_T\).  Hence
\[
  \mathcal N\entails
  \neg\mathsf{Prov}_T(\overline{\ulcorner G_T\urcorner}).
\]
The biconditional (G) is already a theorem of \(\mathsf Q\), not merely an
arbitrary theorem of \(T\); since \(\mathcal N\entails\mathsf Q\), it follows
that \(\mathcal N\entails G_T\).
Therefore \(\neg G_T\) is a false \(\Sigma^0_1\) sentence.  A
\(\Sigma^0_1\)-sound \(T\) cannot prove it, establishing (2).

For (3), suppose \(T\derives\neg G_T\).  By (G),
\[
  T\derives\exists p\,
  \mathsf{Proof}_T(p,\overline{\ulcorner G_T\urcorner}).
\]
For each standard \(n\), consistency and (1) imply that \(n\) is not the code
of a \(T\)-proof of \(G_T\).  Numeralwise representation of the primitive
recursive proof relation gives
\[
  T\derives\neg
  \mathsf{Proof}_T(\bar n,\overline{\ulcorner G_T\urcorner})
\]
for every \(n\).  These statements violate omega-consistency.
The argument is complete once the cited complexity-preserving fixed-point
construction and the cited representability theorem are supplied in full;
those dependencies are why the present status is a proof sketch.
\end{proofplan}

\begin{warning}[What the theorem does not say]
The theorem does not say that every formal theory is deductively incomplete,
that \(G_T\) is unprovable in every stronger theory, or that no deductively
incomplete theory can prove any consistency statement.  The hypotheses and
the subscript \(T\) are essential.
\end{warning}

